\begin{chaptersummary}
  \item A seam is cut by making one service stop claiming what another owns.
        The Sessions service keeps \code{Speaker} as a stub, one field and a
        key declared \code{resolvable: false}, and three things read that
        argument back: \code{implements Node} is gone from the type, the
        \code{\_Entity} union drops it, and the router execution config writes
        \code{disableEntityResolver} beside its key. The flag is advice to the
        composer and not a guard, so the entity route on that service still
        answers a representation for the stub with an unhandled error.
  \item Two services, two seams, two shapes. Speakers takes a whole type away
        and owns it the way any service owns its own; Ratings adds three fields
        to \code{Session} and will never own a session. The second shape is the
        awkward one: it cannot call
        \code{AddGlobalObjectIdentification()}, because that refuses to build a
        schema with no node type in it, and it therefore has to register
        \code{AddDefaultNodeIdSerializer()} by hand in order to read the keys
        it is handed. Leave that line out and the service builds, starts, and
        answers its own root field correctly, while the entity route alone
        fails with \code{Unexpected Execution Error} and nothing in any log.
  \item The first query across two subgraphs costs what it cost inside one.
        Four sessions and three distinct speakers is two statements, one in
        each service, and the batched lookup is
        chapter~\ref{ch:data-without-n-plus-1}'s, reached through a reference
        resolver rather than through a field. The plan grows a second step: a
        \code{BatchEntity} fetch naming the field path it fills in, the fetch it
        waits for, and the representations it will send. The router batches
        those representations no matter how many rows there are, so what the
        service does with a list of them is the subgraph author's problem and
        chapter~\ref{ch:router-n-plus-1}'s subject.
  \item \code{@external} declares a field as somebody else's, and
        \code{@requires} makes the router fetch one before it can answer, even
        when the client never asked for it: \code{feedbackUrl} alone still
        costs the Sessions service a statement for a title nobody sees. The
        field set inside \code{@requires} is checked for being a selection set
        and not for selecting anything, so a misspelled field name reaches the
        composer intact. \code{@provides} is not in this graph, because its two
        preconditions mean keeping a copy of a column in a service that does
        not own it, which is the duplication this chapter's first section
        removed. This composer does not enforce the second of those
        preconditions and Apollo's documentation says it should; write the
        \code{@shareable} anyway, because the stricter of two rules is the one
        that cannot break you later.
  \item Federating on the global node id has a bill and it arrives here.
        \code{@shareable} on \code{node} says any subgraph can answer, and
        neither can: with two node types in two services, \code{node(id:)}
        works when the query names the type in a fragment and fails when it
        names the other one, where a single service answered harmlessly.
        Nothing is broken in either subgraph. The router cannot route on an
        opaque identifier, which is the cost of the opacity
        chapter~\ref{ch:schema-design} was pleased about.
\end{chaptersummary}
